\documentclass[11pt,a4paper]{article}

%% PACHETE MATEMATICE
\usepackage{amsmath,amssymb,amsfonts}
\usepackage{amsthm}
\usepackage{mathpazo}
\usepackage{eulervm}
\usepackage{enumerate}
\usepackage{color}
\usepackage{graphicx}
\usepackage[all]{xy}
\usepackage{xcolor}
\usepackage{framed}
\usepackage[unicode]{hyperref}
\setcounter{MaxMatrixCols}{20}
\usepackage{multicol}
%\usepackage[romanian]{babel}


%% ASPECT
\linespread{1.05}
\usepackage[T1]{fontenc}
\usepackage[utf8x]{inputenc}
\usepackage[margin=2cm]{geometry}
% \usepackage[color]{showkeys}
\usepackage{anyfontsize}
\usepackage{titlesec}
\titleformat*{\section}{\large\bfseries}

\usepackage{fancyhdr}
\pagestyle{fancy}
\rhead{\small \color{gray} Notițe de Adrian Manea}
\lhead{\small \color{gray} BCD-Aero, 2017---2018, L. Costache \& M. Olteanu}


\usepackage[autostyle = false, style = german]{csquotes}
\usepackage[romanian]{babel}
\newcommand\qq{\enquote}

%% COMENZILE MELE
\renewcommand*{\proofname}{Dem.:}
\newcommand\bb{\mathbb}
\newcommand\dr{\mathrm}
\newcommand\kal{\mathcal}
\newcommand\fk{\mathfrak}
\newcommand\op{\oplus}
\newcommand\ot{\otimes}
\newcommand\ds{\displaystyle}
\newcommand\id{\indent}
\newcommand\nid{\noindent}
\newcommand\seq{\subseteq}
\newcommand\sneq{\subsetneq}
\newcommand\speq{\supseteq}
\newcommand\spneq{\supsetneq}
\newcommand\rar{\rightarrow}
\newcommand\Rar{\Rightarrow}
\newcommand\xrar{\xrightarrow}
\newcommand\lar{\leftarrow}
\newcommand\Lar{\Leftarrow}
\newcommand\xlar{\xleftarrow}
\newcommand\lrar{\leftrightarrow}
\newcommand\Lrar{\Leftrightarrow}
\newcommand\ideal{\trianglelefteq}
\newcommand\ai{\text{ a.\^{i}. }}
\newcommand\sk{(\textit{Schi\c{t}\u{a}}) }
\newcommand\rhup{\rightharpoonup}
\newcommand\rhdn{\rightharpoondown}
\newcommand\lhup{\leftharpoonup}
\newcommand\lhdn{\leftharpoondown}
\newcommand\vid{\emptyset}
\newcommand\wh{\widehat}
\newcommand\ol{\overline}
\newcommand\NN{\mathbb{N}}
\newcommand\RR{\mathbb{R}}
\newcommand\ZZ{\mathbb{Z}}
\newcommand\QQ{\mathbb{Q}}
\newcommand\CC{\mathbb{C}}

%% STILURILE MELE
\newtheoremstyle{dotless}{}{}{\itshape}{}{\bfseries}{:}{ }{}

\theoremstyle{dotless}
\newtheorem{theorem}{Teorem\u{a}}[section]
\newtheorem{proposition}{Propozi\c{t}ie}[section]
\newtheorem{lemma}{Lem\u{a}}[section]
\newtheorem{corollary}{Corolar}[section]
\newtheorem{exercise}{Exerci\c{t}iu}

\newtheoremstyle{dotlessdr}{}{}{}{}{\bfseries}{:}{ }{}
\theoremstyle{dotlessdr}
\newtheorem{example}{Exemplu}[section]
\newtheorem{definition}{Defini\c{t}ie}[section]
\newtheorem{remark}{Observa\c{t}ie}[section]




\begin{document}

\begin{center}
  {\large\textbf{Seminar 11 \\
    Probleme de extrem cu forme pătratice}}
\end{center}

\vspace{1cm}


\section{Forme pătratice. Aplicație}
\label{sec:f-patr}

În afară de studiul semnului valorilor proprii ale matricei hessiene, mai putem
decide natura punctelor de extrem pentru o funcție de mai multe variabile și
altfel.

Pentru aceasta, vom utiliza o noțiune algebrică, aceea a \emph{formelor pătratice}.
Mai precis, vom folosi faptul că diferențiala a doua a unei funcții de mai multe
variabile este o formă pătratică.

Din teoria formelor pătratice, ne este de folos următoarea:

\begin{theorem}[Sylvester]\label{thm:sylv-f-patr}
  Fie $\varphi : \RR^n \to \RR, \varphi(x) = \sum a_{ij} x_i x_j$ o formă pătratică
  simetrică ($a_{ij} = a_{ji}, \forall i, j$). Fie matricea sa asociată $A = (a_{ij})$.

  Definim \emph{minorii principali ai matricei} prin:
  \[
    \Delta_1 = a_{11}, %
    \Delta_2 = \begin{vmatrix}
      a_{11} & a_{12} \\
      a_{21} & a_{22}
    \end{vmatrix}, %
    \dots, %
    \Delta_n = \det A.
  \]

  Atunci:
  \begin{enumerate}[(a)]
  \item $\varphi$ este \emph{pozitiv definită} (i.e.\ $\varphi(x) > 0, \forall x$ nenul)
    dacă și numai dacă $\Delta_i > 0, \forall i$;
  \item $\varphi$ este \emph{negativ definită} (i.e.\ $\varphi(x) < 0, \forall x$ nenul)
    dacă și numai dacă $\Delta_1 < 0, \Delta_2 > 0, \dots, (-1)^n \Delta_n > 0$.
  \end{enumerate}
\end{theorem}

Pentru a aplica această teoremă la studiul problemelor de extrem, să remarcăm faptul
că diferențiala a doua a unei funcții de mai multe variabile este o formă pătratică
simetrică. În plus, matricea hessiană este matricea acestei forme pătratice.
Așadar, avem:

\begin{theorem}\label{thm:hes-patr}
  Fie $a \in \RR^n$ un punct critic pentru funcția $f : \RR^n \to \RR$.
  \begin{enumerate}[(a)]
  \item Dacă toți minorii principali ai matricei hessiene sînt strict pozitivi în $a$,
    atunci $a$ este punct de minim local pentru $f$;
  \item Dacă minorii principali ai matricei hessiene au semnele alternate, începînd
    cu primul negativ, atunci $a$ este punct de maxim local;
  \item Dacă toți minorii principali ai matricei hessiene sînt nenuli, dar semnele
    lor nu sînt ca în cazurile de mai sus, atunci $a$ nu este punct de extrem local;
  \item Dacă cel puțin unul din minorii principali ai matricei hessiene este nul,
    nu se poate preciza natura punctului $a$. În acest caz, se evaluează semnul
    diferenței $f(x) - f(a)$, folosind formula lui Taylor pentru funcții de mai
    multe variabile (pînă la termenul de ordin cel puțin 2).
  \end{enumerate}
\end{theorem}

\textbf{Dezvoltarea în serie Taylor pentru o funcție de două variabile reale}
în jurul punctului $a = (x_0, y_0)$ este dată de formula:
\begin{align*}
  f(x, y) &= f(a) + \frac{1}{1!} \Big[ (x - x_0) \frac{\partial f}{\partial x}(a) + %
            (y - y_0) \frac{\partial f}{\partial y}(a) \Big]\\
          &+ \frac{1}{2!} \Big[ (x - x_0)^2 \frac{\partial^2 f}{\partial x^2}(a) + %
            2(x - x_0)(y - y_0) \frac{\partial^2 f}{\partial x \partial y}(a) + %
            (y - y_0)^2 \frac{\partial^2 f}{\partial y^2}(a) \Big] + \\
          &+ \frac{1}{3!} \Big[ (x - x_0)^3 \frac{\partial^3 f}{\partial x^3} (a) + %
            3(x - x_0)^2(y - y_0) \frac{\partial^3 f}{\partial x^2 \partial y}(a) + %
            3(x - x_0)(y - y_0)^3 \frac{\partial^3 f}{\partial x \partial y^2}(a) + %
            (y - y_0)^3 \frac{\partial^3 f}{\partial y^3}(a) \\
          &+ \dots
\end{align*}

\textbf{Exemplu:} Fie funcția
\[
  f : \RR^3 \to \RR, f(x, y, z) = (x - y)^2 + (y - 1)^3 +
  (z - 1)^2.
\]
Găsim punctele critice.

\emph{Soluție:} Rezolvăm sistemul dat de anularea derivatelor parțiale și obținem
punctul $M(1, 1, 1)$. Matricea hessiană are $\Delta_2 = 0$, deci nu putem folosi
teoria formelor pătratice. Folosim definiția.

$M(1, 1, 1)$ este punct de extrem local dacă și numai dacă există o sferă $S((1, 1, 1), r)$
în care diferența $f(x, y, z) - f(1, 1, 1)$ are semn constant (echivalent, într-o dimensiune,
studiem semnul diferenței într-o \emph{vecinătate} a punctului).

În particular, să observăm că $f(x, x, 1) = (x - 1)^3$, care nu are semn constant pentru
$(x, x, 1)$ în orice vecinătate a punctului $(1, 1, 1)$. Deducem că punctul $M$ nu este
de extrem.

%%%%%%%%%%%%%%%%%%%%%%%%%%%%%%%%%%%%%%%%%%%%%%%%%%%%%%%%%%%%%%%%%%%%%%

\section{Exerciții}
\label{sec:exc}

1. Să se arate că funcțiile următoare verifică ecuațiile indicate:
\begin{enumerate}[(a)]
\item $f(x, y) = 2x^2 + 2y^2 + 5xy$,
  $\frac{\partial f}{\partial x} \cdot \frac{\partial f}{\partial y} - %
  5x \frac{\partial f}{\partial x} - 5y \frac{\partial f}{\partial y} + 9xy = 0$;
\item $f(x, y) = xg(x^2 - y^2)$,
  $xy \frac{\partial f}{\partial x} + x^2 \frac{\partial f}{\partial y} = yf$;
\item $f(x, y, z) = g\big(xy, xyz + \frac{x^2}{2} + \frac{x^4}{4} \big)$,
  $-xy \frac{\partial f}{\partial x} + y^2 \frac{\partial f}{\partial y} + %
  x(1 + x^2) \frac{\partial f}{\partial z} = 0$;
\item $f(x, y, z) = g(x - yz, y^2 + z^2)$, $(z^2 - y^2) \frac{\partial f}{\partial x} + %
  z \frac{\partial f}{\partial y} - y \frac{\partial f}{\partial z} = 0$;
\item $f(x, y, z) = g(xyz, x + y + z)$,
  $x(y - z) \frac{\partial f}{\partial x} + y(z - x) \frac{\partial f}{\partial y} +
  z(x - y) \frac{\partial f}{\partial z} = 0$;
\item $f(x, y, z) = g(xyz, yz + zx - xy)$,
  $x^2(y + z) \frac{\partial f}{\partial x} - y^2(z + x) \frac{\partial f}{\partial y} + %
  z^2(y - x) \frac{\partial f}{\partial z} = 0$.
\end{enumerate}

\vspace{1cm}

2. Să se determine punctele de extrem local pentru funcțiile:
\begin{enumerate}[(a)]
\item $f : \RR^2 \to \RR, f(x, y) = x^3 + 4xy - 2x^2 - 2y^2 - 22x + 10 y$;
\item $f : \RR^2 \to \RR, f(x, y) = (x + y^2) \cdot e^{(x + 2y)}$;
\item $f : \RR^2 \to \RR, f(x, y) = e^{(x^2 - y)} + e^{x + y}$;
\item $f : \RR^2 - \{(0, 0)\} \to \RR, f(x, y) = xy \ln (x^2 + y^2)$.
\end{enumerate}

\vspace{1cm}

3. Să se determine punctele de extrem pentru funcțiile $f$ definite pe mulțimile $K$:
\begin{enumerate}[(a)]
\item $f(x, y) = 5x^2 + 4xy + 8y^2$, $K = \{(x, y) \in \RR^2 \mid x^2 + y^2 \leq 1\}$;
\item $f(x, y) = xy^2(x + y - 2)$, $K = \{(x, y) \in \RR^2 \mid x + y \leq 3, x, y \geq 0\}$;
\item $f(x, y) = \frac{x + y}{1 + xy}, \ K = [0, 1] \times [0, 1]$;
\item $f(x, y) = 2x + 3y - \frac{x^2 + y^2}{2}$, $K = \{(x, y) \in \RR^2 \mid x + y \leq 4, %
  5x + 3y \geq 16, x, y \geq 0 \}$;
\item $f(x, y) = x^2 + y^2 + xz - yz$, $K = \{(x, y, z) \in \RR^3 \mid x + y + z = 1, %
  x, y, z \geq 0 \}$.
\end{enumerate}

\end{document}
